\documentclass[11pt]{article}

\usepackage[T1]{fontenc}
\usepackage{lmodern}
\usepackage{microtype}
\usepackage{amsmath,amssymb,amsthm}
\usepackage[margin=1in]{geometry}
\usepackage[hidelinks]{hyperref}

\newtheorem{lemma}{Lemma}
\newtheorem{theorem}[lemma]{Theorem}

\newcommand{\dd}{\mathop{}\!\mathrm{d}}

\title{A Hand-Checkable Proof of Positivity\\
for the Aggregate Curvature Floor}
\author{}
\date{}

\begin{document}
\maketitle

\begin{abstract}
We prove, without numerical computation, the strict positivity of the
one-variable function that supplies the curvature floor in the aggregate-tail
argument for spherical shells. Every estimate is reduced to an exact
rational comparison. In fact, on the full required interval we obtain the
uniform lower bound
\[
 F(\rho)>\frac{1549}{84150}>0.
\]
\end{abstract}

\section{Definitions and elementary estimates}

For \(0\leq s\leq1\), define
\begin{equation}
 G_1(s):=\frac{1}{\pi}
 \left(\sqrt{1-s^2}-s\arccos s\right).
 \label{eq:G1-definition}
\end{equation}
For \(0<\rho<1\), put
\begin{equation}
 \eta_\rho:=G_1\!\left(\max\{\rho,1/2\}\right),
 \qquad
 d_\rho:=\frac{2\arcsin\rho}{\pi},
 \qquad
 b_\rho:=\frac{2\pi\rho}{1-\rho}.
 \label{eq:definitions}
\end{equation}
The function to be estimated is
\begin{equation}
 F(\rho):=2\rho\eta_\rho+d_\rho\rho^2
 -\frac{3(1+d_\rho)b_\rho}{800}.
 \label{eq:F-definition}
\end{equation}
The interval in question is well inside \((0,1)\), since
\[
 0<\frac7{51}<\frac12<\frac{39}{50}<1;
\]
the two nontrivial middle comparisons are respectively \(14<51\) and
\(25<39\). Thus all quantities in \eqref{eq:definitions} and
\eqref{eq:F-definition} are well defined there.

We first record all numerical inequalities used below, together with exact
proofs.

\begin{lemma}[Rational bounds for \(\pi\) and square roots]
\label{lem:rational-bounds}
One has
\begin{equation}
 \pi<\frac{22}{7},\qquad
 \sqrt2>\frac75,\qquad
 \sqrt3>\frac53,\qquad
 \sqrt{\frac{11}{50}}>\frac7{15}.
 \label{eq:rational-bounds}
\end{equation}
\end{lemma}

\begin{proof}
Polynomial division gives the identity
\[
 \frac{x^4(1-x)^4}{1+x^2}
 =x^6-4x^5+5x^4-4x^2+4-\frac4{1+x^2}.
\]
The left-hand side is positive for \(0<x<1\). Hence
\begin{align*}
 0
 &<\int_0^1\frac{x^4(1-x)^4}{1+x^2}\,\dd x\\
 &=\frac17-\frac{4}{6}+1-\frac43+4
   -4\arctan(1)\\
 &=\frac{22}{7}-\pi,
\end{align*}
where \(\arctan(1)=\pi/4\). This proves \(\pi<22/7\).

All three square-root estimates follow by squaring positive quantities:
\[
 2-\left(\frac75\right)^2=\frac1{25}>0,
 \qquad
 3-\left(\frac53\right)^2=\frac29>0,
\]
and
\[
 \frac{11}{50}-\left(\frac7{15}\right)^2
 =\frac{11}{50}-\frac{49}{225}
 =\frac1{450}>0.
\]
\end{proof}

\begin{lemma}[Integral representation and a lower bound for \(G_1\)]
\label{lem:G1-lower}
For \(0\leq s\leq1\),
\begin{equation}
 G_1(s)=\frac1\pi\int_s^1\arccos t\,\dd t
 \geq\frac{2\sqrt2}{3\pi}(1-s)^{3/2}.
 \label{eq:G1-integral-lower}
\end{equation}
In particular, if \(1/2\leq s\leq39/50\), then
\begin{equation}
 G_1(s)>\frac{343}{11250}.
 \label{eq:G1-rational-lower}
\end{equation}
\end{lemma}

\begin{proof}
The antiderivative
\[
 \int\arccos t\,\dd t=t\arccos t-\sqrt{1-t^2}
\]
shows directly that
\[
 \int_s^1\arccos t\,\dd t
 =\sqrt{1-s^2}-s\arccos s,
\]
which proves the first equality in \eqref{eq:G1-integral-lower}.

For \(0\leq t\leq1\), let \(y=\arccos t\), so
\(0\leq y\leq\pi/2\). Since \(\sin u\leq u\) for \(u\geq0\),
\[
 \sqrt{2(1-t)}
 =\sqrt{2(1-\cos y)}
 =2\sin\frac y2
 \leq y=\arccos t.
\]
Integrating this pointwise inequality gives
\begin{align*}
 G_1(s)
 &\geq\frac1\pi\int_s^1\sqrt{2(1-t)}\,\dd t\\
 &=\frac{2\sqrt2}{3\pi}(1-s)^{3/2},
\end{align*}
which completes \eqref{eq:G1-integral-lower}. For completeness,
\(\sin u\leq u\) follows because the derivative of
\(u-\sin u\) is \(1-\cos u\geq0\), and the function vanishes at zero.

Now assume \(1/2\leq s\leq39/50\). Then
\(1-s\geq11/50\). Lemma~\ref{lem:rational-bounds} yields
\begin{align*}
 G_1(s)
 &\geq\frac{2\sqrt2}{3\pi}(1-s)^{3/2}\\
 &\geq\frac{2\sqrt2}{3\pi}
       \left(\frac{11}{50}\right)^{3/2}\\
 &>\frac{2(7/5)}{3(22/7)}
       \frac{11}{50}\frac7{15}\\
 &=\frac{49}{165}\frac{77}{750}
  =\frac{343}{11250}.
\end{align*}
Here the last equality follows by cancelling a factor \(11\) from
\(3773/123750\).
\end{proof}

\section{Strict positivity of the curvature floor}

\begin{theorem}
\label{thm:F-positive}
For every
\[
 \frac7{51}\leq\rho\leq\frac{39}{50},
\]
the function in \eqref{eq:F-definition} satisfies
\begin{equation}
 F(\rho)>\frac{1549}{84150}>0.
 \label{eq:uniform-F-lower}
\end{equation}
\end{theorem}

\begin{proof}
Because \(\rho>0\), substitution of
\(b_\rho=2\pi\rho/(1-\rho)\) into \eqref{eq:F-definition} and division
by \(\rho\) give the exact identity
\begin{equation}
 \frac{F(\rho)}{\rho}
 =2\eta_\rho+d_\rho\rho
 -\frac{3\pi(1+d_\rho)}{400(1-\rho)}.
 \label{eq:F-over-rho}
\end{equation}
We treat the two branches in the definition of \(\eta_\rho\) separately.

\medskip
\noindent
\textbf{First branch: \(7/51\leq\rho\leq1/2\).}
Here \(\max\{\rho,1/2\}=1/2\), and therefore
\begin{align}
 \eta_\rho
 &=G_1(1/2)
 =\frac1\pi\left(\frac{\sqrt3}{2}
 -\frac12\arccos\frac12\right)\notag\\
 &=\frac{\sqrt3}{2\pi}-\frac16,
 \label{eq:eta-low-branch}
\end{align}
because \(\arccos(1/2)=\pi/3\). Since \(\arcsin\) is increasing on
\([0,1]\),
\[
 0<d_\rho=\frac{2\arcsin\rho}{\pi}
 \leq\frac{2\arcsin(1/2)}{\pi}=\frac13.
\]
Also \(1-\rho\geq1/2\), so \((1-\rho)^{-1}\leq2\). It follows that
\begin{equation}
 \frac{3\pi(1+d_\rho)}{400(1-\rho)}
 \leq\frac{3\pi(4/3)\,2}{400}
 =\frac\pi{50}
 <\frac{11}{175}.
 \label{eq:low-negative-bound}
\end{equation}
On the other hand, Lemma~\ref{lem:rational-bounds} gives
\begin{equation}
 2\eta_\rho
 =\frac{\sqrt3}{\pi}-\frac13
 >\frac{5/3}{22/7}-\frac13
 =\frac{35}{66}-\frac{22}{66}
 =\frac{13}{66}.
 \label{eq:low-eta-bound}
\end{equation}
The middle term \(d_\rho\rho\) in \eqref{eq:F-over-rho} is nonnegative.
Combining \eqref{eq:F-over-rho}, \eqref{eq:low-negative-bound}, and
\eqref{eq:low-eta-bound}, we obtain
\begin{equation}
 \frac{F(\rho)}{\rho}
 >\frac{13}{66}-\frac{11}{175}
 =\frac{2275-726}{11550}
 =\frac{1549}{11550}.
 \label{eq:low-F-over-rho}
\end{equation}
Since \(\rho\geq7/51\) and \(11550=7\cdot1650\), this implies
\begin{equation}
 F(\rho)>\frac7{51}\frac{1549}{11550}
 =\frac{1549}{84150}.
 \label{eq:low-uniform-F}
\end{equation}

\medskip
\noindent
\textbf{Second branch: \(1/2\leq\rho\leq39/50\).}
Here \(\max\{\rho,1/2\}=\rho\), so
\(\eta_\rho=G_1(\rho)\). Lemma~\ref{lem:G1-lower} gives
\begin{equation}
 2\eta_\rho>\frac{343}{5625}.
 \label{eq:high-eta-bound}
\end{equation}
Monotonicity of \(\arcsin\), together with \(\rho\geq1/2\), gives
\begin{equation}
 d_\rho\geq\frac13,
 \qquad
 d_\rho\rho\geq\frac16.
 \label{eq:high-d-lower}
\end{equation}
For the required upper bound on \(d_\rho\), observe that
\[
 \left(\frac{39}{50}\right)^2
 =\frac{1521}{2500}
 <\frac{1875}{2500}=\frac34
 =\left(\frac{\sqrt3}{2}\right)^2.
\]
Both sides are positive, hence \(39/50<\sqrt3/2\). Consequently
\begin{equation}
 d_\rho
 \leq\frac{2\arcsin(39/50)}{\pi}
 <\frac{2\arcsin(\sqrt3/2)}{\pi}
 =\frac23.
 \label{eq:high-d-upper}
\end{equation}
Finally, \(1-\rho\geq11/50\), so
\((1-\rho)^{-1}\leq50/11\). From
\eqref{eq:high-d-upper} and Lemma~\ref{lem:rational-bounds},
\begin{align}
 \frac{3\pi(1+d_\rho)}{400(1-\rho)}
 &<\frac{3\pi(5/3)}{400}\frac{50}{11}
 =\frac{5\pi}{88}\notag\\
 &<\frac{5(22/7)}{88}
 =\frac5{28}.
 \label{eq:high-negative-bound}
\end{align}
Substitution of \eqref{eq:high-eta-bound},
\eqref{eq:high-d-lower}, and \eqref{eq:high-negative-bound} into
\eqref{eq:F-over-rho} yields
\begin{align}
 \frac{F(\rho)}{\rho}
 &>\frac{343}{5625}+\frac16-\frac5{28}\notag\\
 &=\frac{9604+26250-28125}{157500}\notag\\
 &=\frac{7729}{157500}.
 \label{eq:high-F-over-rho}
\end{align}
Since \(\rho\geq1/2\),
\begin{equation}
 F(\rho)>\frac{7729}{315000}.
 \label{eq:high-uniform-F}
\end{equation}
The high-branch lower bound is stronger than the one asserted in
\eqref{eq:uniform-F-lower}, because exact subtraction gives
\[
 \frac{7729}{315000}-\frac{1549}{84150}
 =\frac{120341}{19635000}>0.
\]

The two branches cover the full interval. Their definitions agree at
\(\rho=1/2\), because
\[
 G_1\!\left(\max\{1/2,1/2\}\right)=G_1(1/2).
\]
Consequently, \eqref{eq:low-uniform-F} and
\eqref{eq:high-uniform-F} prove \eqref{eq:uniform-F-lower} everywhere
on \([7/51,39/50]\).
\end{proof}

\end{document}
